%======================================================================================
\section{Conclusions and limits}\label{sec:conclusions}

\begin{minipage}[t]{0.49\textwidth}
\textbf{Supported}
\begin{itemize}
  \item groups are mostly pairs; large ones are strings along the route that change members every few minutes,
        while the episode lasts about an hour;
  \item group pilots keep the same pace: their delay wanders half as fast as that of two pilots kept in
        contact by chance, and they stay closer after the contact;
  \item the size distribution follows from one pilot joining or leaving at a time; its heavy tail from members
        of larger groups leaving less often;
  \item the lead changes a few times per hour and goes more often to the experienced pilot on the higher wing;
  \item pilots return to the same partners and to their club mates more than who flies where and when explains.
\end{itemize}
\end{minipage}\hfill
\begin{minipage}[t]{0.49\textwidth}
\textbf{Not supported}
\begin{itemize}
  \item a spring that pulls pilots back: the pull is about that of chance contacts selected the same way;
  \item that the shape of a group (triangle or chain, complete or tree) changes how its pilots fly;
  \item a social rank from the order in a group: a line along the route is geometry;
  \item that keeping pace means following: the surrogate pairs fly at another hour, so pilots in the same
        air at the same time may share the pace without seeing each other;
  \item pilots are names of the catalog: homonyms join, a pilot who changes name splits.
\end{itemize}
\end{minipage}
